\color{unidarkgreen}{\subsection{Контроль синхронизма и напряжений (КСН)}}
\color{black}

\begin{enumerate}[label=\arabic{section}.\arabic{subsection}.\arabic{enumi}, labelsep=4pt, leftmargin=0pt, itemindent=57pt, itemsep=0pt, parsep=0pt] % если что можно крутить itemsep

\item
Функциональный блок контроля синхронизма и напряжений (КСН) предназначен для формирования
разрешающих сигналов, соответствующих условиям включения, для АПВ и оперативного включения.

\item
Величина напряжения ошиновки для КСН (\textit{«UТ(АТ)»}) поступает от трехфазного ТН и используется без
дополнительной коррекции. Величины напряжений энергообъектов (\textit{«Uэ1»} и \textit{«Uэ2»}) могут быть получены как
от трехфазного ТН, так и от однофазного ТН или ШОН. Тип ТН на энергообъекте выбирается уставкой \textbf{«Тип ТН Эn»}.

\item
Выбор напряжения (фазного или линейного), используемого для КСН, выполняется единой для всех
энергообъектов уставкой \textbf{«Фаза подключения КСН»}. Если на энергообъекте установлен трехфазный ТН, то
напряжение от него передается в КСН без дополнительной корректировки, и уставка \textbf{«Фаза подключения
КСН»} определяет только тип используемого напряжения. Если же применяется однофазный ТН или ШОН, то
напряжение отбора физически должно быть подключено к входу напряжения фазы А устройства. В этом случае
уставка \textbf{«Фаза подключения КСН»} должна быть установлена в то положение, которое соответствует фазе
фактического подключения однофазного ТН или ШОН.

\item
Учет разности напряжений между энергообъектом и ошиновкой, возникающий при использовании
различных ТН или особенностях работы ШОН, выполняется при подключении однофазного ТН или ШОН. Для
этого в устройстве предусмотрена возможность корректировки напряжения энергообъекта с помощью уставок
\textbf{«K подстройки Uэn»} и \textbf{«φ подстройки Uэn»}. Корректировка по амплитуде и по фазе осуществляется по
следующим выражениям:

\begin{align}
U_{\textmd{Эn}} = K_{\textmd{под Эn}} \cdot U_{\textmd{отб Эn}}, \\
φ_{U_{\textmd{Эn}}} = φ_{U_{\textmd{отб Эn}}} + φ_{\textmd{под Эn}}, 
\end{align}
\begin{eqexpl} [10mm]
    \item{$ U_{\textmd{отб Эn}} $} действующие значения напряжения отбора Эn, подключенное к аналоговому входу напряжения фазы А, В;
	\item{$ U_{\textmd{под Эn}} $} значение уставки \textbf{«K подстройки Uэn»};
	\item{$ φ_{U_{\textmd{отб Эn}}} $} фаза напряжения отбора Эn, эл. градусы;
	\item{$ φ_{\textmd{под Эn}} $} значение уставки \textbf{«φ подстройки Uэn»}, эл. градусы;
	\item{$ φ_{U_{\textmd{Эn}}} $} фаза напряжения отбора Эn с учетом корректировки, эл. градусы.
\end{eqexpl}


\begin{figure}[H]
\centering
\includegraphics[width=0.75\textwidth,height=0.75\textheight,keepaspectratio]{\fbpath/761.3351 КСН/img/U_KSN.png}
\captionof{figure}{Выбор величин напряжений для КСН}\label{pict_ksn:U_KSN}
\end{figure}

\item
Основные режимы включения и контролируемые для них условия приведены в таблице \ref{ksn:tbl1}. Режимы 
включения с контролем синхронизма (КС) и с улавливанием синхронизма (УС) вводятся соответствующим
 программными ключами \textbf{«КС Вn»} и \textbf{«УС Вn»}.

\item
Если разность частот напряжений шин и линии не превышает уставку \textbf{«df макс КС»}, то
разрешается включение с контролем синхронизма. Для режима включения с КС также контролируется:
\begin{itemize}
    \item наличие напряжения на обоих энергообъектах (ошиновка и линия);
    \item разность напряжений ошиновки и линии, не превышающая по модулю значения заданной уставки;
    \item разность фаз напряжений ошиновки и линии, не превышающая по модулю значения заданной уставки.
\end{itemize}

\item
При выводе программного ключа \textbf{«КС Вn»} включение разрешается без контроля частоты, а только по
уровню напряжений и разности их фаз.

\item
Если разность частот превышает уставку \textbf{«df макс КС Вn»}, но не выходит за пределы уставки \textbf{«df пред.
доп. Вn»}, то разрешается включение с улавливанием синхронизма. Сигнал включения с улавливанием
синхронизма формируется с опережением на время включения выключателя. Режим включения с УС можно
вывести программным ключом \textbf{«УС Вn»}.

\item
Наличие или отсутствие напряжения на линии контролируется измерительными органами
максимального и минимального напряжения, пороги срабатывания которых задаются уставками \textbf{«Uнал Т(АТ)
Bn»} и \textbf{«Uотс Эn»} соответственно. Пороги срабатывания измерительных органов для контроля
напряжения на энергообъекте задаются уставками \textbf{«Uнал Эn»} и \textbf{«Uотс Эn»}.

\small
\setlength{\extrarowheight}{0.05cm} % высота ячеек в таблице
\begin{longtable}
{|>{\centering\arraybackslash}m{2.7cm}|>{\raggedright\arraybackslash}m{14.2cm}|}
\caption{Условия включения\hfill\vspace{-0.5\baselineskip}}\label{ksn:tbl1} \\ 
\hhline{|-|-|}
\rowcolor{gray!30}
Режим включения & Условия включения  \\ 
\hhline{|-|-|}
\endfirsthead

\caption*{\hspace{3pt}\emph{Продолжение таблицы \ref{ksn:tbl1}\hfill\vspace{-0.5\baselineskip}}} \\ % сделано по ГОСТ 2.105 п.6.8.7
\hhline{|-|-|}
\rowcolor{gray!30}
Режим включения & Условия включения  \\ 
\hhline{|-|-|}
\endhead
\hhline{|-|-|}

\endfoot
\endlastfoot

ОНТ(АТ)+ННЭn &
\makecell[lt]{%
– $U_{\textsl{Т(АТ)}}$ меньше уставки \textbf{«Uотс Т(АТ) Вn»}; \\
– $U_{\textsl{Эn}}$ больше уставки \textbf{«Uнал Эn»}. } \\
\hline

ОНЭn+ННТ(АТ) &
\makecell[lt]{%
– $U_{\textsl{Т(АТ)}}$ больше уставки \textbf{«Uнал Т(АТ) Вn»}; \\
– $U_{\textsl{Эn}}$ меньше уставки \textbf{«Uотс Эn»}. } \\
\hline

КС(УС) &
\makecell[lt]{%
1) Включение контролем синхронизма: \\
 – $U_{\textsl{Т(АТ)}}$ больше уставки \textbf{«Uнал Т(АТ) Вn»}; \\
 – $U_{\textsl{Эn}}$ больше уставки \textbf{«Uнал Эn»}; \\
 – разность напряжений ошиновки и энергообъекта Эn не превышает по модулю \\ значения заданной уставки \textbf{«dU макс КС(УС) Вn»}; \\
 – разность фаз напряжений ошиновки и энергообъекта Эn не превышает по модулю \\ значения заданной уставки \textbf{«dФ макс КС Вn»}; \\
 – разность частот напряжений ошиновки и энергообъекта Эn не превышает значения \\ заданной уставки \textbf{«df макс КС Вn»}; \\
2) Включение с улавливанием синхронизма: \\
 – $U_{\textsl{Т(АТ)}}$ больше уставки \textbf{«Uнал ЛЭП Вn»}; \\
 – $U_{\textsl{Э1}}$ больше уставки \textbf{«Uнал Эn»};\\
 – разность напряжений ошиновки и энергообъекта Эn не превышает по модулю \\ значения заданной уставки \textbf{«dU макс КС(УС) Вn»}; \\
 – разность частот напряжений ошиновки и энергообъекта Эn превышает значение \\ уставки \textbf{«df макс КС Вn»}; \\
 – разность частот напряжений ошиновки и энергообъекта Эn не превышает значения \\ заданной уставки \textbf{«df пред. доп. Вn»}; \\
 – за время \textbf{«Т опер. вкл. Вn»} с момента подачи сигнала на включение выключателя \\ расхождение углов между векторами напряжений на ошиновке и энергообъекте Эn \\ не превысит $10^\circ$.} \\
\hline

\end{longtable}
\setlength{\extrarowheight}{0.0cm} % возврат параметра высота ячеек в таблице
\normalsize

\item
Сигнал разрешения АПВ \textit{«Разр. АПВ Вn»} формируется при выполнении условий включения в режимах
ОНТ(АТ)+ННЭn и КС(УС), либо ОНЭn+ННТ(АТ) и КС(УС.

\item
Сигнал разрешения оперативного включения \textit{«Разр. опер. вкл. Вn»} может формироваться от внутренней
логики КСН или от внешнего сигнала в зависимости от положения программного переключателя
\textbf{«Контр. опер. вкл. Вn»}.

\item
В положении \textit{«Внутр.»} оперативное включение разрешается при выполнении условий включения в
режимах ОНТ(АТ)+ННЭn (при опробовании ошиновки), ОНЭn+ННТ(АТ) (при опробовании энергообъекта), КС
или УС (при замыкании в транзит). Также предусмотрена возможность оперативного включения с контролем
отсутствия напряжения на ошиновке или на энергообъекте при введенном программном ключе \textbf{«Опер. вкл. Вn
с КОН»}. При отсутствии внешнего сигнала \textit{«Опер. вкл. с КС(УС) Вn»} оперативное включение может
выполняться без контроля напряжений и условий синхронизма.

\item
В положении \textit{«Внеш.»} оперативное включение разрешается при наличии внешнего сигнала
\textit{«Внеш. сигнал КС(УС) Вn»}.

\begin{figure}[H]
\centering
\includegraphics[width=0.85\textwidth,height=0.85\textheight,keepaspectratio]{\fbpath/761.3351 КСН/img/KSN.pdf}
\captionof{figure}{Функциональный блок КСН В1}\label{pict_ksn:KSN}
\end{figure}


\small
\begin{longtable}{|>{\centering\arraybackslash}m{5.3cm}|>{\centering\arraybackslash}m{3.3cm}|>{\centering\arraybackslash}m{4.2cm}|>{\centering\arraybackslash}m{1.8cm}|>{\centering\arraybackslash}m{1cm}|}
\caption{Параметры для настройки КСН\hfill\vspace{-0.5\baselineskip}}\label{ksn:tbl1}\\ 
\hline
\rowcolor{gray!30}
Параметр & Условное обозначение на схеме & Значение/ Диапазон & Единица измерения & Шаг \\ 
\hline
\endfirsthead
\caption*{\hspace{3pt}\emph{Продолжение таблицы \ref{ksn:tbl1}\hfill\vspace{-0.5\baselineskip}}} \\ % сделано по ГОСТ 2.105 п.6.8.7
\hline
\rowcolor{gray!30}
Параметр & Условное обозначение на схеме & Значение/ Диапазон & Единица измерения & Шаг \\ 
%\hline
\endhead

%===m>КСН|-

%===m
\end{longtable}
\normalsize

\end{enumerate}